\section{Conclusion}

\looseness=-1
This paper turns a common concern about safety alignment into a concrete
measurement problem. In the Gemma 4 pairs, the aligned models usually keep
working but produce less grounded artifacts, while the less-restricted
derivatives also do better on non-security controls. In Qwen2.5-Coder and
Llama 3.1, the same pattern does not hold. The result is therefore not a clean
security-specific alignment effect and not a general endorsement of
less-restricted derivatives. It is evidence that autonomous security-agent
behavior depends on grounding, capability, derivative provenance, and
tool-interface reliability.

The practical conclusion is that safety evaluations for security agents should
be trace-based and evidence-aware. They should ask not only whether a model
refuses, but whether it reads the right local evidence, uses authorized tools
without escaping the sandbox, names the relevant program objects, and produces
a final artifact that a defender could act on. They should also test ordinary
coding controls, because a security-task gap is hard to interpret without a
baseline for general agent reliability. In this study, removing alignment
recovers some grounded behavior for Gemma, but it introduces unsafe-action
events for the 26B derivative and does not solve hard proof and patch
verification. The right object of
evaluation is the whole agentic system: the model, tools, authorization
boundary, trace evidence, and final checker together. That is the standard a
security-agent benchmark should measure, not a transcript-level refusal rate
alone.
